\documentclass[11pt,a4paper]{article}
\usepackage[T1]{fontenc}
\usepackage{lmodern}
\usepackage[margin=23mm]{geometry}
\usepackage{amsmath,amssymb,bm}
\usepackage[hidelinks]{hyperref}
\usepackage{fancyhdr}
\pagestyle{fancy}
\fancyhf{}
\fancyhead[L]{Phase--amplitude equations}
\fancyhead[R]{Transcribed notes}
\fancyfoot[C]{\thepage}
\setlength{\headheight}{14pt}
\setlength{\parindent}{0pt}
\setlength{\parskip}{5pt}
\newcommand{\dd}{\mathrm{d}}
\newcommand{\er}{\hat{\bm e}_r}
\newcommand{\ez}{\hat{\bm z}}
\newcommand{\ex}{\hat{\bm x}}
\newcommand{\ey}{\hat{\bm y}}
\begin{document}
\begin{center}
{\LARGE\bfseries Phase--Amplitude Equations}\\[5pt]
{\large Organized transcription of four handwritten pages}
\end{center}

\textit{Editorial convention.} A dot denotes $\dd/\dd\tau$. Derivative
primes in the original are replaced by dots; parameter primes, such as $k'$
and $\gamma'$, are retained. Crossed-out intermediate text is omitted.
Some symbol readings are provisional; see the transcription notes on page~3.

\section{Angular and radial equations}
For the first page, write $\bm r=u\ex+v\ey$, with
$u=r\cos\theta$, $v=r\sin\theta$, and $r^2=u^2+v^2$.
The angular derivative is
\begin{equation}
\dot\theta
=\frac{\dfrac{\dd}{\dd\tau}(v/u)}{1+(v/u)^2}
=\frac{u^2\dfrac{\dot v u-\dot u v}{u^2}}{u^2+v^2}
=\frac{\dot v u-\dot u v}{u^2+v^2}.
\end{equation}
The transverse vector equation in the notes reads
\begin{align}
\frac{\dd}{\dd\tau}(u\ex+v\ey)
&=-\left[(u\ex+v\ey)\times\ez
+k'(\chi-1)\ez\times\bm b'\right] \notag\\
&=-\left[-u\ey+v\ex
+k'(\chi-1)(\cos\theta_B\,\ey-\sin\theta_B\,\ex)\right].
\end{align}
Thus
\begin{equation}
\boxed{\begin{aligned}
\dot u&=-v+k'(\chi-1)\sin\theta_B,\\
\dot v&=u-k'(\chi-1)\cos\theta_B.
\end{aligned}}
\end{equation}
With $\delta=\theta_B-\theta$,
\begin{align}
\dot v u-\dot u v
&=v^2+u^2-k'(\chi-1)r
 (\sin\theta_B\sin\theta+\cos\theta_B\cos\theta)\notag\\
&=r^2-k'(\chi-1)r\cos\delta,
\end{align}
so that
\begin{equation}
\boxed{\dot\theta=1-\frac{k'(\chi-1)\cos\delta}{r}.}
\end{equation}
The handwritten reference ``from Eq.~(23)'' gives
\begin{align}
\dot r
&=\er\cdot\frac{\dd\bm r}{\dd\tau}
=-\er\cdot\left[\bm r\times\ez
+k'(\chi-1)\ez\times\bm b'\right]\notag\\
&=-k'(\chi-1)(\ez\times\bm b')\cdot\er\notag\\
&=-k'(\chi-1)(\bm b'\times\er)\cdot\ez\notag\\
&=k'(\chi-1)\sin\delta.
\end{align}

\newpage
\section{Relative phase and coupled equations}
The wave-phase derivative is transcribed as follows. The middle factors
in this particular line are less clear in the photograph.
\begin{align}
\dot\theta_B
&=\frac{\dd}{\dd\tau}(k'z')
=-k'\beta_z'\frac{\gamma'}{\Omega_\lambda}c\notag\\
&=-\frac{k}{\gamma_T}\beta_z'\frac{\gamma'}{\Omega_\lambda}c
=-\frac{\omega}{c}\frac{m\gamma'}{\Omega_\lambda\gamma_T}c\beta_z'
\notag\\
&=-2\gamma'\beta_z'=-(\chi-1).
\end{align}
Consequently,
\begin{align}
\dot\delta&=\dot\theta_B-\dot\theta\notag\\
&=-(\chi-1)-1+\frac{k'(\chi-1)\cos\delta}{r}\notag\\
&=-\chi+\frac{k'(\chi-1)\cos\delta}{r}.
\end{align}
The three equations collected at the top of the second page are
\begin{equation}
\boxed{\begin{aligned}
\dot\delta&=-\chi+\frac{k'(\chi-1)\cos\delta}{r},\\
\dot r&=k'(\chi-1)\sin\delta,\\
\dot\chi&=-2\gamma'\varepsilon-rk'\sin\delta.
\end{aligned}}
\label{eq:system}
\end{equation}

\section{Eliminating the explicit phase term}
The stated objective is to cancel the phase term on the right-hand side.
First,
\begin{align}
\dot\chi-r\dot r
&=-2\gamma'\varepsilon-rk'\sin\delta
-rk'\chi\sin\delta+rk'\sin\delta\notag\\
&=-2\gamma'\varepsilon-rk'\chi\sin\delta.
\end{align}
The notes propose adding $\dd(k'r\cos\delta)/\dd\tau$.
Indeed,
\begin{align}
\frac{\dd}{\dd\tau}(r\cos\delta)
&=\dot r\cos\delta-r\sin\delta\,\dot\delta\notag\\
&=k'(\chi-1)\sin\delta\cos\delta
-r\sin\delta\left(-\chi+\frac{k'(\chi-1)\cos\delta}{r}\right)
\notag\\
&=k'(\chi-1)\sin\delta\cos\delta+r\chi\sin\delta
-k'(\chi-1)\sin\delta\cos\delta\notag\\
&=\chi r\sin\delta.
\end{align}
For constant $k'$, this yields
\begin{equation}
\dot\chi-r\dot r+\frac{\dd}{\dd\tau}(k'r\cos\delta)
=-2\gamma'\varepsilon,
\end{equation}
or equivalently
\begin{equation}
\boxed{\frac{\dd}{\dd\tau}
\left(\chi-\frac{r^2}{2}+k'r\cos\delta\right)
=-2\gamma'\varepsilon.}
\end{equation}

\newpage
\section{Phase-based Cartesian variables and $J$}
The second page now introduces
\begin{equation}
u=r\cos\delta,\qquad v=r\sin\delta.
\label{eq:phasecoords}
\end{equation}
These reuse the letters $u,v$ for components defined by $\delta$, whereas
Section~1 used components defined by $\theta$.
Substitution gives
\begin{equation}
\frac{\dd}{\dd\tau}
\left(\chi-\frac{u^2+v^2}{2}+k'u\right)
=-2\gamma'\varepsilon.
\end{equation}
Introduce
\begin{equation}
\boxed{J=\chi-\frac{u^2+v^2}{2}+k'u
=\chi-\frac{(u-k')^2+v^2}{2}+\frac{k'^2}{2}.}
\end{equation}
The handwritten integration is
\begin{equation}
\frac{\dd J}{\dd\tau}=-2\gamma'\varepsilon
\quad\Longrightarrow\quad
\boxed{J=J_0-2\gamma'\varepsilon\tau.}
\end{equation}
This displayed linear form presumes constant $\gamma'\varepsilon$ and
$J_0=J(0)$.

Finally, differentiating the variables in Eq.~\eqref{eq:phasecoords} gives
\begin{align}
\frac{\dd u}{\dd\tau}
&=\dot r\cos\delta-r\sin\delta\,\dot\delta
=\chi r\sin\delta=\chi v,\\
\frac{\dd v}{\dd\tau}
&=\dot r\sin\delta+r\cos\delta\,\dot\delta\notag\\
&=k'(\chi-1)-u\chi
=-\chi(u-k')-k'.
\end{align}
Hence the final pair on the second page is
\begin{equation}
\boxed{\dot u=\chi v,\qquad \dot v=-\chi(u-k')-k'.}
\end{equation}

\section*{Transcription notes for the first two pages}
\begin{enumerate}
\item The handwritten variable resembling $x$ is typeset as $\chi$.
The primed factor in $-2\gamma'\varepsilon$ is read as $\gamma'$;
its handwriting also resembles $\nu'$. These are provisional symbol readings.
\item In the expanded $\dot\theta_B$ line, $\Omega_\lambda$, $m$, and
$\gamma_T$ are provisional readings of the handwriting. Their definitions
are not supplied in these photographs. This line is transcribed, not
independently validated.
\item Overwritten signs in $\dot v u-\dot u v$ and $\dot\theta$
are rendered as minus signs, consistent with the immediately preceding
component equations and the subsequent $\dot\delta$ equation.
\item The original reference to Eq.~(23) is preserved as a reference in the
handwritten source; that equation itself was not supplied separately.
\item The two definitions of $u,v$ are intentionally retained and identified.
The assumption that $k'$ is constant, and the condition for linear $J(\tau)$,
are editorial clarifications. No definitions of the physical parameters
have been added.
\item Partially visible writing on adjacent pages, crossed-out text, and
unrelated marginal fragments are excluded.
\end{enumerate}
\newpage
\section{Shifted variables and the complex equation}
The third handwritten page continues with the shifted variables
\begin{equation}
q=u-k',\qquad p=v.
\end{equation}
For constant $k'$, the final equations of Section~4 become
\begin{equation}
\boxed{\dot q=\chi p,\qquad \dot p=-\chi q-k'.}
\end{equation}
Set $z=q+\mathrm{i}p$ and combine the two real equations:
\begin{align}
\dot q+\mathrm{i}\dot p
&=\chi p-\mathrm{i}\chi q-\mathrm{i}k',\notag\\
\dot z&=-\mathrm{i}\chi(q+\mathrm{i}p)-\mathrm{i}k'
=-\mathrm{i}\chi z-\mathrm{i}k'.
\label{eq:complex}
\end{align}

\section{Introducing the detuning parameter $\Lambda$}
Using the previous expression for $J$,
\begin{equation}
J=\chi-\frac{q^2+p^2}{2}+\frac{k'^2}{2},
\qquad \Lambda=\frac{q^2+p^2}{2}-\chi.
\end{equation}
It follows that
\begin{align}
\Lambda
&=\frac{k'^2}{2}-J
=\frac{k'^2}{2}-(J_0-2\gamma'\varepsilon\tau)\notag\\
&=\frac{k'^2}{2}-J_0+2\gamma'\varepsilon\tau
=\Lambda_0+2\gamma'\varepsilon\tau,
\qquad \Lambda_0=\frac{k'^2}{2}-J_0.
\end{align}
Therefore $\chi=-\Lambda+(q^2+p^2)/2=-\Lambda+|z|^2/2$.
Substitution into Eq.~\eqref{eq:complex} yields
\begin{equation}
\dot z=-\mathrm{i}\left(\frac{|z|^2}{2}-\Lambda\right)z
-\mathrm{i}k'
\quad\Longrightarrow\quad
\boxed{\mathrm{i}\dot z-\left(\frac{|z|^2}{2}-\Lambda\right)z=k'.}
\end{equation}

\section{Comparison with a nonlinear Schr\"odinger-type equation}
The reference form written in the notes is
\begin{equation}
\mathrm{i}\frac{\dd\Psi}{\dd\tau}+(|\Psi|^2-\tau)\Psi=\mu,
\qquad \mu\in\mathbb{R}.
\end{equation}
Taking the complex conjugate gives
\begin{equation}
-\mathrm{i}\frac{\dd\Psi^*}{\dd\tau}
+(|\Psi|^2-\tau)\Psi^*=\mu.
\end{equation}
For $z=-\Psi^*$, the last handwritten line reads
\begin{equation}
\frac{\dd z}{\dd\tau}+(\tau-|\Psi|^2)z=\mu
\qquad\text{(as written in the photograph).}
\end{equation}
\textit{Editorial correction:} direct substitution retains a factor
$\mathrm{i}$ in front of the derivative. Thus the corrected line is
\begin{equation}
\boxed{\mathrm{i}\frac{\dd z}{\dd\tau}+(\tau-|\Psi|^2)z=\mu,
\qquad |z|^2=|\Psi|^2.}
\end{equation}
\textit{Scope of the comparison:} the two complex equations still have
different coefficients ($|z|^2/2$ versus $|z|^2$, and $\Lambda(\tau)$ versus
$\tau$). The additional handwritten page supplies the scaling in
Section~8 below. As before, linear $\Lambda(\tau)$ assumes
constant $k'$ and $\gamma'\varepsilon$.
\newpage
\section{Normalization and the threshold written in the notes}
The fourth handwritten page starts from
\begin{equation}
\mathrm{i}\frac{\dd z}{\dd\tau}
-\left(\frac{|z|^2}{2}-\Lambda\right)z=k',
\qquad \Lambda=\Lambda_0+2\gamma'\varepsilon\tau.
\end{equation}
The proposed substitutions are $z=AZ$, $\tau=BT$, and $\Lambda=CT$.
With constant, real $A$ and $B$, these give
\begin{equation}
\mathrm{i}\frac{\dd Z}{\dd T}
-Z\left(\frac{A^2B}{2}|Z|^2-BCT\right)=\frac{k'B}{A}.
\end{equation}
Choose
\begin{equation}
\frac{A^2B}{2}=1,\qquad BC=1,
\qquad a=2\gamma'\varepsilon>0,\qquad C=\sqrt a.
\end{equation}
Then the scaling factors in the notes are
\begin{equation}
B=\frac{1}{\sqrt a},\qquad
A=(2\sqrt a)^{1/2}=\sqrt{2}\,a^{1/4},\qquad
\frac{A}{B}=\sqrt{2}\,a^{3/4}.
\end{equation}
The dimensionless forcing and the resulting equation are
\begin{equation}
\mu=\frac{k'}{A/B}=\frac{k'}{\sqrt{2}\,a^{3/4}},
\qquad
\boxed{\mathrm{i}\frac{\dd Z}{\dd T}-(|Z|^2-T)Z=\mu.}
\end{equation}

\textit{Time-origin clarification.} The simultaneous choices $\tau=BT$
and $\Lambda=CT$ require $\Lambda_0=0$. For nonzero $\Lambda_0$, retain
the constant through a shift of the time origin:
\begin{equation}
\boxed{T=\sqrt a\left(\tau+\frac{\Lambda_0}{a}\right),\qquad
\tau=\frac{T}{\sqrt a}-\frac{\Lambda_0}{a},\qquad
\Lambda=\sqrt a\,T.}
\end{equation}
This shift leaves $\dd T/\dd\tau=\sqrt a$ unchanged, so the amplitude
scaling and the normalized equation above still apply.

The last two lines of the photograph state the autoresonance threshold as
\begin{equation}
\mu>0.41
\quad\Longrightarrow\quad
k'>0.41\sqrt{2}\,a^{3/4}\simeq0.58\,a^{3/4}.
\end{equation}
Using the additional handwritten identification $a=\xi_E/(cB_0)$,
this is written as
\begin{equation}
\boxed{k'\gtrsim0.58\left(\frac{\xi_E}{cB_0}\right)^{3/4}.}
\end{equation}
\textit{Transcription note.} The numerator in $\xi_E/(cB_0)$ is read
provisionally as $\xi_E$; its definition is not given on this page.
The numerical threshold $0.41$ is reproduced from the handwriting.
The photograph does not specify its initial-condition assumptions or
provide a derivation of that numerical value.
\end{document}
